\TieudeT{PHONG TỎA VẬT LÍ}
\subsubsection*{PHẦN I. Thí sinh trả lời câu hỏi từ câu 1 đến câu 18. Mỗi câu hỏi thí sinh chỉ chọn một phương án}
\setcounter{ex}{0}
\begin{ex}%book%[1P2Y1-2][Loai1]
Hai lực $\vec F_1$ và $\vec F_2$ cùng phương, ngược chiều, cùng điểm đặt và có độ lớn $|F_1| = |F_2| = 50\ N$. Hợp lực của chúng là lực nào sau đây?
\choice
{\True $F = 0\ N$}
{$F = 100\ N$; cùng phương cùng chiều với lực $\vec F_1$}
{$F = 100\ N$; cùng phương cùng chiều với lực $\vec F_2$}
{$F = 25\ N$; cùng phương cùng chiều với lực $\vec F_1$}
\haidong{1}
\loigiai{ 
Do hai lực có độ lớn bằng nhau và cùng phương, ngược chiều $\Rightarrow$ hợp lực bằng 0.
}
%<MyLT>
\haidong{7} \end{ex}
\begin{ex}
Cho hai lực đồng quy có độ lớn lần lượt là $16\ N$ và $12\ N$. Độ lớn hợp lực của chúng \textbf{không} thể là
\choice
{$F=20\ N$}
{\True $F=30\ N$}
{$F=5\ N$}
{$F=15\ N$}
\loigiai{}
\end{ex}
\begin{ex}%book%[1P2Y1-2][Loai1]
Cho hai lực đồng quy có độ lớn $F_1 = F_2 = 60\ N$. Góc tạo bởi hai lực là $120^\circ$. Độ lớn của hợp lực là
\choice
{$60\sqrt2\ N$}
{$30\ N$}
{$60\sqrt3\ N$}
{\True $60\ N$}
\haidong{3}
\loigiai{Độ lớn của hợp lực là $F = 2F_1\cos {60^\circ} = F_1 = 60$ N}
\haidong{7} \end{ex}
\begin{ex}%book%[1P2Y1-2][Loai1]
Cho hai lực đồng quy có độ lớn $F_1=16\ N$; $F_2=12\ N$. Cho độ lớn của hợp lực là $20\ N$. Góc giữa hai lực là
\choice
{$0^\circ$}
{$60^\circ$}
{\True $90^\circ$}
{$120^\circ$}
\haidong{3}
\loigiai{ 
\begin{itemize}
\item Ta nhận thấy do $F^2 = F_1^2 + F_2^2 $
\item Nên theo định lí Pitago $\Rightarrow$ $\alpha = 90^\circ$.
\end{itemize}
}
\haidong{7} \end{ex}
\begin{ex}%book %[1P2B1-2][Loai1]
Cho hai lực đồng quy có độ lớn $F_1=F_2=20\ N$. Độ lớn của hợp lực là $F=34,6\ N$ khi đó hai lực thành phần hợp với nhau một góc là
\choice
{$30^\circ $}
{\True $60^\circ $}
{$90^\circ $}
{$120^\circ $}
\haidong{3}
\loigiai{
Góc giữa hai lực là $60^\circ$.
}

\haidong{7} %<MyLT1>
\end{ex}
\begin{ex}%book
Một chất điểm chịu tác dụng đồng thời của hai lực thành phần vuông góc với nhau có độ lớn lần lượt là $F_{1}=10\ N$ và $F_{2}$. Biết hợp lực trên có độ lớn là $10 \sqrt{5}\ N$. Giá trị của $F_{2}$ là
\choice
{$10\ N$}
{\True $20\ N$}
{$30\ N$}
{$40\ N$}

\haidong{4}
\loigiai{Chọn B

}\haidong{7} %<MyLT1>
\end{ex}
\begin{ex}%book
Cho ba lực đồng phẳng, đồng quy $\vec{F}_{1}$, $\vec{F}_{2}$, $\vec{F}_{3}$ có độ lớn lần lượt là $16\ N$, $12\ N$ và $12\ N$. Biết góc tạo bởi các lực $\left(\vec{F}_{1},\ \vec{F}_{2}\right)=30^{\circ}$ và $\left(\vec{F}_{2},\ \vec{F}_{3}\right)=120^{\circ}$. Độ lớn hợp lực của ba lực này bằng
\choice
{$27,62\ N$}
{$31,7\ N$}
{$16\ N$}
{\True $20\ N$}

\haidong{3}
\loigiai{Chọn D

}\haidong{7} \end{ex}
\begin{ex}%book %[1P2K1-2][Loai1]
Một vật chịu bốn lực tác dụng. Lực $F_1 = 40\ N$ hướng về phía Đông, lực $F_2 = 50\ N$ hướng về phía Bắc, lực $F_3 = 70\ N$ hướng về phía Tây và lực $F_4 = 90\ N$ hướng về phía Nam. Độ lớn của hợp lực tác dụng lên vật là
\choice
{\True $50\ N$}
{$131\ N$}
{$170\ N$}
{$50\ N$}
\haidong{3}
\loigiai{ 
\begin{itemize}
\item Hợp lực của hai lực hướng theo hướng Đông - Tây: $F_{13} = F_3 - F_1 = 30$ N (do $F_3> F_1$ và các lực cùng phương; ngược chiều)
\item Tương tự ta có hợp lực của hai lực hướng theo hướng Nam - Bắc: $F_{24} = F_4 - F_2 = 40$ N (do $F_4> F_2$ và các lực cùng phương; ngược chiều)
\item Vì các hướng Đông - Tây và Nam - Bắc vuông góc nhau $\Rightarrow$ hợp lực cần tính.
\item $F = \sqrt {F_{13}^2 + F_{24}^2 } = \sqrt {30^2 + 40^2 } = 50\ N$.
\end{itemize}	
}
\haidong{7} \end{ex}



\begin{ex}%book%[1P2B1-3][Loai1]
Một chất điểm đứng yên dưới tác dụng của ba lực $12\ N$, $15\ N$, $9\ N$. Góc giữa hai lực có độ lớn $12\ N$ và $9\ N $ bằng
\choice
{$30^\circ$}
{$60^\circ$}
{\True $90^\circ$}
{$120^\circ$}
\haidong{3}
\loigiai{
Độ lớn ba lực tuân theo bộ số Pitago nên góc giữa hai lực nhỏ là $90^\circ$.
}
\haidong{8} %<MyLT1>
\end{ex}
\begin{ex}%book
Một vật rắn cân bằng dưới tác dụng của ba lực có độ lớn $F_{1}=6\ N$, $F_{2}=8\ N$, $F_{3}=10\ N$. Nếu bỏ đi lực $F_{2}$ thì hợp lực của hai lực còn lại có độ lớn là
\choice
{$10\ N$}
{\True $8\ N$}
{$16\ N$}
{$14\ N$}	
\haidong{2}
\loigiai{Chọn B

}\haidong{8} \end{ex}
\begin{ex}%book%[1P2-5.4K][Loai1]
Một người nặng $75\ kg$, khi ở trên mặt trăng có $g = 1,7\ m/s^2$ thì trọng lượng của người đó xấp xỉ bằng
\choice
{$750\ N$}
{\True $127\ N$}
{$254\ N$}
{$735\ N$}
\loigiai{
Trọng lượng của người đó khi ở trên Mặt Trăng là: $P = mg = 75.1,7 = 127,5$ N.	
}
\haidong{10} \end{ex}
\begin{ex}%book %[1P2B5-4][Loai1]
 Trọng lượng một người có khối lượng $75\ kg$, tại nơi có gia tốc trọng trường $g = 9,80\ m/s^2$ là 
\choice
{\True $735\ N$}
{$750\ N$}
{$75\ N$}
{$0\ N$}
\loigiai{ 
Ở Trái Đất : $P_{TD} = m.g_{TD} = 75.9,8 = 735$ N.
}
%<MyLT>
\haidong{10} \end{ex}	
\begin{ex}%book
Gia tốc rơi tự do trên bề mặt Trái Đất, trên bề mặt Mặt Trăng và trên bề mặt Kim Tinh lần lượt là $9,8 \ m / s^2, 1,6 \ m / s^2$ và $8,7 \ m / s^2$. Trọng lượng của một nhà du hành vũ trụ có khối lượng 75 kg khi người đó ở trên Trái Đất, trên Mặt Trăng và trên Kim Tinh lần lượt là $P_1, P_2$ và $P_3$. Độ lớn của $\left(P_1+P_2-P_3\right)$ bằng
\choice
{$179 \ N$}
{$205 \ N$}
{\True $203 \ N$}
{$275 \ N$}
\loigiai{
$
\begin{aligned}
& \text {Chọn C} \\
& \left(P_1+P_2-P_3\right)=75 \cdot(9,8+1,6-8,7)=202,5 \approx 203 N
\end{aligned}
$

}
\haidong{10} \end{ex}	
\begin{ex}%book%SBT KNTT 17.4
Một người đi chợ dùng lực kế để kiểm tra khối lượng của một gói hàng. Người đó treo gói hàng vào lực kế và đọc được số chỉ của lực kế là $20\ N$. Biết gia tốc rơi tự do tại vị trí này là $g=10\ m/s^2$. Khối lượng của túi hàng là
\choice
{\True $2\ kg$}
{$20\ kg$}
{$30\ kg$}
{$10\ kg$}
\loigiai{
\begin{itemize}
\item Ta có: $P=mg\Rightarrow m=\dfrac{P}{g}=\dfrac{20}{10}=2\ kg$.
\end{itemize}
}
\haidong{10} \end{ex}	
\begin{ex}%book%SBT KNTT 17.8
Một vật nặng có khối lượng $0,2\ kg$ được treo vào một sợi dây không dãn. Lấy $g=9,8\ m/s^2$. Lực căng của dây khi vật cân bằng là
\choice
{$9,8\ N$}
{\True $1,96\ N$}
{$0,2\ N$}
{$2,25\ N$}
\loigiai{
\begin{itemize}
\item Khi quả cầu treo dưới sợi dây cân bằng thì lực căng sợi dây bằng trọng lượng của vật: $T=P=mg=0,2.9,8=1,96\ N$.
\end{itemize}
}
\haidong{10} \end{ex}	
\begin{ex}%book
Một vật có khối lượng 10 kg được treo bằng hai sợi dây. Mỗi sợi dây tạo  một góc $45^\circ$ so với phương ngang. Lấy $g=10 \ m / s^2$. Độ lớn của lực căng dây trong mỗi sợi dây bằng
\choice
{\True $70,7 \ N$}
{$92,3 \ N$}
{$100,4 \ N$}
{$60,1 \ N$}
\loigiai{
Chọn A
Áp dụng ĐKCB $\vec{T_1}+\vec{T_2}+\vec{P}=0$

Theo phương thẳng đứng: $P=T_1 \sin \alpha_1+T_2 \sin \alpha_2$

Mà $T_1=T_2=T$ và $\alpha_1=\alpha_2=\alpha=45^\circ$

$$
P=2 T \sin \alpha \Rightarrow T=\dfrac{P}{2 \sin \alpha}=70,7 N
$$

}
\haidong{10} \end{ex}	
\begin{ex}%book
Một khối $30\ kg$ được giữ trên mặt phẳng nghiêng $30^\circ$ không ma sát bằng một sợi dây gắn cố định vào tường sao cho dây luôn song song với mặt phẳng nghiêng. Lấy $g=10\ m/s^2$. Lực căng dây có độ lớn là
\choice
{\True $150 \ N$}
{$150 \sqrt{3} \ N$}
{$150 \sqrt{2} \ N$}
{$300 \ N$}
\loigiai{
Chọn A
$T=P \sin \alpha=mg \sin \alpha=150 \ N$

}
\haidong{10} \end{ex}	
\begin{ex}%Tailieuchuan
\immini[thm]{Một vật khối lượng $m=5,0 \ kg$ nằm yên trên một mặt phẳng nghiêng nhờ một sợi dây song song với mặt phẳng nghiêng. Góc nghiêng $\alpha=30^\circ$. Lấy $g=10 \ m/s^2$. Lực căng của dây có độ lớn là
\choice
{\True $T=25\ N$}
{$T=50\ N$}
{$T=43\ N$}
{$T=25\ N$}}{
\begin{tikzpicture}[scale=1,thick,>=stealth',draw=Mapcolor]
\tkzInit[ymin=-4,ymax=0.5,xmin=-0.5,xmax=5.5]\tkzClip
\begin{scope}[shift={(0,0)},rotate=-30]
\node  at  (2.3,0)  [shape=rectangle,draw=red,fill=green!70!black,minimum height=0.5cm,minimum width=0.6cm, rounded corners=1pt,rotate=-30]  {};
\draw[very thick,->,Mapcolor](2.3,0)--++(0:0.7)node[above,black,xshift=10pt]{}coordinate(px);
\draw[very thick,->,Mapcolor](2.3,0)--++(-180:0.9)node[above,black]{$\vec T$}coordinate(n);
\draw[very thick,->,Mapcolor](2.3,0)--++(-60:1.4)node[right,black,yshift=5pt]{$\vec{P}$}coordinate(p);
\draw[dashed,thin](px)--(p);
\draw[fill=white](2.3,0)circle(1.2pt);
\draw(0,-0.27)coordinate(a)--++(0:6)coordinate(b)--++(210:5)coordinate(c);
\fill[pattern=north east lines](b)--(c)--++(300:0.1)--++(390:5);
\end{scope}
\draw (a) -- (0.08,0.13);

\draw (0.05,-0.02) -- (1.68,-0.95);
\draw[very thick,->,Mapcolor] (2.02,-1.14) -- (2.4,-0.5);
\node at (2.6,-0.4) {$\vec N$};
\end{tikzpicture}
}
\loigiai{}
\end{ex}











\subsubsection*{PHẦN 2. Thí sinh trả lời từ câu 1 đến câu 4. Trong mỗi ý a), b), c), d) ở mỗi câu, thí sinh chọn đúng hoặc sai.}
\setcounter{ex}{0}
\begin{ex}
\immini[thm]{Một chất điểm chịu tác dụng của hai lực $\vec{F_1}$, $\vec{F_2}$, hợp với nhau một góc $\alpha$. Độ lớn của hợp lực là $F$.}{	\begin{tikzpicture}[scale=1,thick,>=stealth',draw=Mapcolor]
\coordinate (o) at (0,0);
\draw[->,very thick,Mapcolor](o)--++(0:3.5)coordinate(f2)node[below]{\normalsize{\bth{\vec{F}_1}}};
\draw[->,very thick,Mapcolor](o)--++(60:2)coordinate(f1)node[above]{\normalsize{\bth{\vec{F}_2}}};
\draw[dashed,Mapcolor](f1)--++(0:3.5)coordinate(f) node[right]{\normalsize{\bth{\vec{F}}}};
\draw[black,dashed] (f)--(f2);
\draw[->,very thick,Mapcolor](o)--++(f);
\end{tikzpicture}}
\choiceTF[t]
{$F^2=F_1^2+F_2^2-2 F_1 F_2 \cos \alpha$}
{\True $F_{\max}=F_1+F_2$}
{$F_{\min}=0$}
{\True $\left|F_1-F_2\right| \leq F \leq F_1+F_2$}
\loigiai{
\begin{itemchoice}
\itemch
\includegraphics[width=4cm]{images/2025_06_04_af03f1d88c157fe9bf98g-105(1)}
$
F^2=F_1^2+F_2^2+2 F_1 F_2 \cos \alpha
$

\itemch $F_{\text {max}}=F_1+F_2$.
$\Rightarrow$ Đúng.
\itemch $F_{\text {min}}=\left|F_1-F_2\right|$.
$\Rightarrow$ Sai
\itemch $\left|F_1-F_2\right| \leq F \leq F_1+F_2$.
$\Rightarrow$ Đúng.
\end{itemchoice}

}
\end{ex}	
\begin{ex}%Tailieuchuan
Cho hai lực đồng quy có độ lớn bằng nhau $F_1=F_2=20 \ N$.
\choiceTF[t]
{\True Độ lớn của hợp lực khi chúng hợp với nhau một góc $\alpha=0^\circ$ là $40\ N$} 
{Độ lớn của hợp lực khi chúng hợp với nhau một góc $\alpha=90^\circ$ là $20 \sqrt{6} \ N$} 
{\True Độ lớn của hợp lực khi chúng hợp với nhau một góc $\alpha=60^\circ$ là $20 \sqrt{3}\ N$} 
{\True Độ lớn của hợp lực là $20\ N$ khi chúng hợp với nhau một góc $\alpha=120^\circ$} 
\loigiai{}
\end{ex}	
\begin{ex}
Một bóng đèn được treo tại chính giữa một dây nằm ngang làm dây bị võng xuống. Biết trọng lượng của đèn là $100\ N$ và góc giữa hai nhánh dây là $120^{\circ}$. 
\choiceTF[t]
{Lực căng của hai nhánh dây cùng phương ngược chiều với nhau}
{\True Lực căng của hai nhánh dây có độ lớn bằng nhau}
{\True Lực tổng hợp từ lực căng của hai nhánh dây cùng phương, ngược chiều với trọng lực tác dụng lên bóng đèn}
{Lực căng của mỗi nhánh dây là $50\ N$}
\loigiai{\begin{itemchoice}
\itemch Sai.
\itemch Đúng. Điểm treo bóng đèn nằm chính giữa dây, do đó: $T_{1}=T_{2}=T$
\itemch Đúng. Vì bóng đèn đang nằm cân bằng nên: $\vec{T}_{1}+\vec{T}_{2}+\vec{P}=\vec{0} \Rightarrow \vec{T}_{1}+\vec{T}_{2}=-\vec{P}$
\itemch Sai. Độ hợp lực của $\vec{T}_{1}$ và $\vec{T}_{2}$ là: $T_{12}=2 T \cos \dfrac{120^{\circ}}{2}=2 T \cos 60^{\circ}$\\ Từ điều kiện cân bằng, ta có: $T_{12}=P=2 T \cos 60^{\circ} \Rightarrow T=\dfrac{P}{2 \cos 60^{\circ}}=\dfrac{100}{2 \cos 60^{\circ}}=100 ~N$
\end{itemchoice}
}

\haidong{10} \end{ex}
\begin{ex}
\immini[thm]{Người ta treo một cái đèn trọng lượng $P = 5\ N$ vào một giá đỡ gồm hai thanh cứng nhẹ $AB$ và $AC$ như hình vẽ. Biết $\alpha = 60^\circ$, $AB$ vuông góc với $BC$, lấy $g = 10\ m/s^2$. Gọi $\vec{T}$ và $\vec{N}$ lần lượt là lực căng của thanh $AC$ và phản lực của tường tác dụng lên thanh $AB$. Biết phản lực của tường lên thanh $AB$ có phương trùng với phương của thanh.
}{\begin{tikzpicture}[scale=1,thick,>=stealth',draw=Mapcolor]
\tkzInit[ymin=-4.5,ymax=0.5,xmin=-0.5,xmax=3]\tkzClip
\tkzDefPoints{0/0/b, 0/0.5/m, 0/-3/n, -0.1/-3/p, -0.1/0.5/q}
\tkzDefShiftPoint[b](-45:3){a}
\tkzDefShiftPoint[a](-90:1){v}
\tkzDefPointBy[projection=onto m--n](a)\tkzGetPoint{o}
\fill[pattern=north east lines](m)--(n)--(p)--(q);
\tkzDrawSegments[line width=3pt,gray](o,a)
\tkzDrawSegments(a,b m,n a,v)
\tkzDrawPoints[size=2,fill=white](a)
\tkzMarkAngles[size=0.7cm](o,b,a)
\tkzLabelAngles[pos=1.2](o,b,a){$\alpha$}
\tkzLabelPoint[left](b){$C$}
\tkzLabelPoint[right](a){$A$}
\draw[fill=Mapcolor!30] (v) ellipse (0.2cm and 0.3cm);

\draw[fill=Mapcolor!30]  (2,-2.6) rectangle (2.25,-2.89);
\node at (-0.3,-2.1) {$B$};
\end{tikzpicture}}

\choiceTF[t]
{$\vec{T} + \vec{N} = \vec{P}$}
{\True $T = 10\ N$}
{$N = P = 5\ N$}
{\True Lực mà thanh $AB$ tác dụng vào cột $BC$ gần bằng $8,66\ N$}
\loigiai{
\begin{itemchoice}
\itemch Khi đèn cân bằng:

\includegraphics[width=4cm]{images/2025_06_04_af03f1d88c157fe9bf98g-108(2)}


$$
\vec{P}+\vec{N}+\vec{T}=\vec{0} \rightarrow \vec{N}+\vec{T}=-\vec{P}
$$

$\Rightarrow$ Sai
\itemch $T=\dfrac{P}{\cos \alpha}=10N$.
$\Rightarrow$ Đúng

\itemch $N=\sqrt{T^2-P^2}=5 \sqrt{3}N$.
$\Rightarrow$ Sai
\itemch Theo định luật III Newton thì lực mà thanh $A B$ tác dụng lên cột $BC$ có độ lớn bằng lực mà cột $BC$ tác dụng lên $A B$ bằng $5 \sqrt{3} N \approx 8,66N$.
\end{itemchoice}
}
\end{ex}

\subsubsection*{PHẦN 3. Thí sinh trả lời từ câu 1 đến câu 6.}
\setcounter{ex}{0} 
\begin{ex}
Một chất điểm chịu tác dụng của hai lực $F_1 = 9\ N$ và $F_2 = 12\ N$, hợp với nhau một góc $0^\circ$. Độ lớn của hợp bằng bao nhiêu niutơn?
\shortans[oly]{21}
\loigiai{
$\alpha = 0^\circ \Rightarrow F = F_1 + F_2 = 21\ N$
}
\end{ex}
\begin{ex}
Một chất điểm chịu tác dụng của hai lực $F_1 = 9\ N$ và $F_2 = 12\ N$, hợp với nhau một góc $180^\circ$. Độ lớn của hợp bằng bao nhiêu niutơn?
\shortans[oly]{3}
\loigiai{
$\alpha = 180^\circ \Rightarrow F = \left| F_1 - F_2 \right| = 3\ N$
}
\end{ex}
\begin{ex}
Một chất điểm chịu tác dụng của hai lực $F_1=10\ N$ và $F_2=24\ N$, hợp với nhau một góc $90^\circ$. Độ lớn của hợp bằng bao nhiêu niutơn?
\shortans[oly]{26}
\loigiai{
$\alpha=90^\circ \rightarrow F=\sqrt{F_1^2+F_2^2}=26 N$.
}
\end{ex}
\begin{ex}
\immini[thm]{Một thùng gỗ khối lượng 10 kg được treo vào một sợi dây nằm cân bằng trên mặt phẳng nghiêng $30^\circ$ so với phương ngang. Lấy $g=9,8\ m/s^2$. Lực căng dây là bao nhiêu niutơn?}{\begin{tikzpicture}[scale=1,thick,>=stealth',draw=Mapcolor]
\tkzInit[ymin=-4,ymax=0.5,xmin=-0.5,xmax=5.5]\tkzClip
\begin{scope}[shift={(0,0)},rotate=-30]
\node  at  (2.3,0)  [shape=rectangle,draw=red,fill=green!70!black,minimum height=0.5cm,minimum width=0.6cm, rounded corners=1pt,rotate=-30]  {};

\draw(0,-0.27)coordinate(a)--++(0:6)coordinate(b)--++(210:5)coordinate(c);
\fill[pattern=north east lines](b)--(c)--++(300:0.1)--++(390:5);
\end{scope}
\draw (a) -- (0.08,0.13);

\draw (0.02,-0.01) -- (1.72,-1);

\node at (4.1,-3) {$30^\circ$};
\end{tikzpicture}}
\shortans[oly]{49}
\loigiai{
49N
}
\end{ex}
\begin{ex}%DoanhNamDinh
Một nhà du hành vũ trụ khi ở mặt đất thì có trọng lượng là $784\ N$, khi ở bề mặt của Kim Tinh thì có trọng lượng là $696\ N$. Biết rằng gia tốc rơi tự do trên bề mặt Trái Đất là $9,8 \ m/s^2$. Gia tốc rơi tự do ở bề mặt Kim Tinh có độ lớn là bao nhiêu $m/s^2$?
\shortans[oly]{8,7}
\loigiai{

\begin{itemize}
\item Vận dụng công thức $P=m\ g$. Ta có $m=\dfrac{P_{T D}}{g_{T D}}=\dfrac{P_{K T}}{g_{K T}} \Rightarrow g_{K T}=\dfrac{P_{K T}}{P_{T D}} g_{T D}=8,7 \ m/s^2$
\end{itemize}

}
\end{ex}
\begin{ex}%DoanhNamDinh
\immini[thm]{Một thùng hàng trọng lượng $500\ N$ đang trượt xuống dốc. Mặt dốc tạo với phương ngang một góc $30^\circ$. Áp lực do thùng hàng tác dụng lên mặt phẳng nghiêng có độ lớn bằng bao nhiêu niutơn (làm tròn kết quả đến chữ số hàng đơn vị)?}{\includegraphics[width=4cm]{images/2025_06_04_af03f1d88c157fe9bf98g-180(1)}}
\shortans[oly]{433}
\loigiai{
\begin{itemize}
\item Áp lực do thùng hàng tác dụng lên mpn là $\vec{P_{y}}$

\includegraphics[width=4cm]{images/2025_06_04_af03f1d88c157fe9bf98g-180}

\includegraphics[width=4cm]{images/2025_06_04_af03f1d88c157fe9bf98g-181(1)}

\item Ta có $P_{y}=P \cdot \cos 30^\circ=433,01\ N$.
\item Kết quả lấy phần nguyên $\rightarrow P_{y}=433\ N$
\end{itemize}
}
\end{ex}





\newpage 




\TieudeT{PHONG TỎA VẬT LÍ}
\subsubsection*{PHẦN I. Thí sinh trả lời câu hỏi từ câu 1 đến câu 18. Mỗi câu hỏi thí sinh chỉ chọn một phương án}
\setcounter{ex}{0}
\begin{ex}
Hai lực thành phần $\vec{F_1}$ và $\vec{F_2}$ có độ lớn lần lượt là $F_1$ và $F_2$, hợp lực $\vec{F}$ của chúng có độ lớn là $F$. Ta có
\choice
{$F$ luôn lớn hơn $F_1$}
{$F$ luôn nhỏ hơn $F_2$}
{\True $\left|F_1-F_2\right| \leq F \leq F_1+F_2$}
{$F$ không thể bằng $F_1$}
\loigiai{}
\end{ex}
\begin{ex}%book%[1P2-1.2Y][Loai1]
Cho hai lực đồng quy có độ lớn bằng $9\ N$ và $12\ N$ hợp với nhau góc $90^\circ$. Độ lớn của hợp lực là
\choice
{$1\ N$}
{$3\ N$}
{\True $15\ N$}
{$21\ N$}
\haidong{3}
\loigiai{      
Độ lớn của hợp lực: $F_{hl} = \sqrt {F_1^2 + F_2^2 } = \sqrt {9^2 + 12^2 } = 15\ N.$
}
\haidong{7} \end{ex}
\begin{ex}%book
Một chất điểm chịu tác dụng của đồng thời hai lực cùng độ lớn $20\ N$, góc hợp bởi hai lực là $120^{\circ}$. Hợp lực của hai lực trên
\choice
{\True $F=20\ N ;\ \left(\vec{F} ;\ \vec{F}_{1}\right)=60^{\circ}$}
{$F=20 \sqrt{3}\ N ;\ \left(\vec{F} ;\ \vec{F}_{1}\right)=60^{\circ}$}
{$F=20\ N ;\left(\vec{F} ;\ \vec{F_{1}}\right)=120^{\circ}$}
{$F=20 \sqrt{3}\ N ;\ \left(\vec{F} ;\ \vec{F_{1}}\right)=120^{\circ}$}
\haidong{5}
\loigiai{Chọn A}\haidong{7} \end{ex}
\begin{ex}%book%[1P2-1.2B][Loai1]
Cho hai lực đồng quy có cùng độ lớn $100\ N$. Góc giữa hai lực bằng bao nhiêu thì hợp lực cũng có độ lớn bằng $100\ N$?
\choice
{\True $120\degree$}
{$90\degree$}
{$180\degree$}
{$0\degree$}
\haidong{3}
\loigiai{
Hai lực đồng quy có cùng độ lớn $\Rightarrow F = 2F_1\cos \dfrac{\alpha}{2} \Rightarrow \dfrac{\alpha}{2} = \dfrac{F}{2F_1} = \dfrac{100}{2.100} = 0,5  \Rightarrow \alpha = {120^\circ}$.
}
\haidong{7} \end{ex}
\begin{ex}%book
Hai lực $\vec{F}_{1}$ và $\vec{F}_{2}$ vuông góc với nhau. Có độ lớn là $3\ N$ và $4\ N$. Hợp lực của chúng tạo với hai lực này các góc lần lượt là 
\choice
{$30^{\circ}$ và $60^{\circ}$}
{$42^{\circ}$ và $48^{\circ}$}
{\True $37^{\circ}$ và $53^{\circ}$}
{$45^{\circ}$ và $45^{\circ}$}

\haidong{3}
\loigiai{Chọn C

}\haidong{7} \end{ex}
\begin{ex}%book
Một chất điểm chịu tác dụng của ba lực đồng phẳng, đồng quy có cùng độ lớn $15\ N$. Biết góc tạo bởi các lực $\left(\vec{F}_{1}, \vec{F}_{2}\right)=\left(\vec{F}_{2}, \vec{F}_{3}\right)=60^{\circ}$. Độ lớn hợp lực của ba lực bằng
\choice
{\True $30\ N$}
{$20\ N$}
{$15\ N$}
{$45\ N$}

\haidong{5}
\loigiai{Chọn A

}\haidong{7} \end{ex}
\begin{ex}%book%[1P2-1.2B][Loai1] 
\immini[thm]{Cho bốn lực đồng quy, đồng phẳng như hình vẽ bên. Biết $F_1=5\ N$, $F_2=3\ N$, $F_3=7\ N$, $F_{4}=1\ N$. Hợp lực của bốn lực đó có độ lớn là
\choice
{\True $2\sqrt{2}\ N$}
{$2\ N$}
{$3\ N$}
{$3\sqrt{3}\ N$}
}{\begin{tikzpicture}[scale=0.5,thick,>=stealth',draw=Mapcolor]
\tkzDefPoints{0/0/o, -5/0/1, 0/3/2, 7/0/3, 0/-1/4}
\tkzDrawSegments[very thick,Mapcolor,->](o,1 o,2 o,3 o,4)
\tkzDrawPoints[size=3,fill=black](o)
\tkzLabelPoint[above](1){$\vec{F}_1$}
\tkzLabelPoint[right](2){$\vec{F}_2$}
\tkzLabelPoint[above](3){$\vec{F}_3$}
\tkzLabelPoint[right](4){$\vec{F}_4$}
\end{tikzpicture}}
\haidong{4}
\loigiai{
\begin{itemize}
\item Ta có: $\vec{F}=\vec{F}_1+\vec{F}_2+\vec{F}_3+\vec{F}_{4}=\left(\vec{F}_2+\vec{F}_{4}\right)+\left(\vec{F}_1+\vec{F}_3\right)=\vec{F}_{24}+\vec{F}_{13}$.  
$\\
\text{Với}\,\vec{F}_{24}:
\begin{cases}
\vec{F}_{24}\uparrow \uparrow \vec{F}_2\\
F_{24}=F_2-F_{4}=2\ N\\
\end{cases}
$\\
$
\text{Với}\,\vec{F}_{13}:
\begin{cases}
\vec{F}_{13}\uparrow \uparrow \vec{F}_3\\
F_{13}=F_3-F_1=2\ N\\
\end{cases}
$
\item Dễ suy ra $\vec{F}_{24}\bot \vec{F}_{13}\Rightarrow F=\sqrt{F_{13}^2+F_{24}^2}=2\sqrt{2}\ N$.
\end{itemize}
}
\haidong{7} \end{ex}

\begin{ex}%book%[1P2K1-3][Loai1]
\immini[thm]{Một vật chịu tác dụng của ba lực ${\vec{F}}_{1}$, $\vec{F}_2$, $\vec{F}_3$ như hình vẽ bên thì nằm cân bằng. Biết rằng độ lớn của lực $F_3 = 40\sqrt{3}\ N$. Độ lớn của lực $F_1$ và $F_2$ bằng
\choice
{$F_1=60\ N$ và $F_2=80\ N$}
{\True $F_1=40\ N$ và $F_2=80\ N$}
{$F_1=50\ N$ và $F_2=60\ N$}
{$F_1=40\ N$ và $F_2=60\ N$}
}{\begin{tikzpicture}[scale=0.5,thick,>=stealth',draw=Mapcolor]
\tkzDefPoints{0/0/o, 1.5/0/a, -0.5/-1/i}
\tkzDefShiftPoint[o](120:3){b}
\tkzDefPointBy[translation=from o to b](a) \tkzGetPoint{c}
\tkzDefPointBy[symmetry=center o](c)\tkzGetPoint{d}
\tkzDrawSegments[very thick,Mapcolor,->](o,a o,b o,d)
\tkzDrawPoints[size=3,fill=black](o)
\tkzMarkAngles[size=0.4cm](a,o,b)
\tkzLabelAngle[pos=0.9](a,o,b){$120^{\circ}$}
\tkzLabelPoint[below](a){$\vec{F}_1$}
\tkzLabelPoint[above](b){$\vec{F}_2$}
\tkzLabelPoint[below left](d){$\vec{F}_{3}$}
\end{tikzpicture}} 
\haidong{3}
\loigiai{
\immini{\begin{itemize}
\item Để vật cân bằng thì lực tổng hợp của hai lực $\vec{F}_1$ và $\vec{F}_2$ phải cùng phương, ngược chiều với $\vec{F}_2$. 
\item Gọi $\vec{F}_{13}=\vec{F}_1+\vec{F}_3\Rightarrow \vec{F}_{13}$ phải tạo với $\vec{F}_1$ một góc $60^\circ$. 
\item Từ hình vẽ có: $\cos 30^\circ=\dfrac{F_3}{F_{13}}\Rightarrow F_{13}=\dfrac{F_3}{\cos 30^\circ}=80\ N\Rightarrow F_2=80\ N$
\item Vì $\vec{F}_1\bot \vec{F}_3\Rightarrow F_1=\sqrt{F_{13}^2-F_3^2}=40\ N$.
\end{itemize}}{\begin{tikzpicture}[scale=0.7,thick,>=stealth']
\tkzDefPoints{0/0/o, 1.5/0/a, -0.5/-1/i}
\tkzDefShiftPoint[o](-90:2.6){b}
\tkzDefPointBy[translation=from o to b](a) \tkzGetPoint{c}
\tkzDefPointBy[symmetry=center o](c)\tkzGetPoint{d}
\tkzDrawSegments[very thick,blue,->](o,a o,b o,d)
\tkzDrawSegments[very thick,red,->](o,c)
\tkzDrawSegments[dashed](a,c b,c)
\tkzDrawPoints[size=3,fill=black](o)
\tkzMarkAngles[size=0.4cm](a,o,d)
\tkzLabelAngle[pos=0.9](a,o,d){$120^{\circ}$}
\tkzLabelPoint[above](a){$\vec{F}_1$}
\tkzLabelPoint[left](b){$\vec{F}_3$}
\tkzLabelPoint[above right](c){$\vec{F}_{13}$}
\tkzLabelPoint[left](d){$\vec{F}_2$}
\end{tikzpicture}}
}
\haidong{8} \end{ex}
\begin{ex}%book%[1P2B1-3][Loai1]
Một chất điểm đứng yên dưới tác dụng của ba lực $8\ N$, $10\ N$, $12\ N$. Nếu bỏ đi lực $10\ N$ thì hợp lực của hai lực còn lại là
\choice
{$20\ N $}
{$4\ N $}
{$6\ N $}
{\True $10\ N $}
\haidong{3}
\loigiai{
Để vật cân bằng thì hợp của hai lực còn lại phải cân bằng với lực đầu tiên nên có giá trị là 10 N.
}
%<MyLT>
\haidong{8} \end{ex}
\begin{ex}%book
Một vật có khối lượng $m=6\ kg$ được giữ yên trên một mặt phẵng nghiêng bởi một sợi dây song song với đường dốc chính. Độ cao và chiều dài dốc lần lượt là $5\ m$ và $10\ m$; lấy $g=10\ m / s^{2}$ và ma sát không đáng kể. Lực căng của sợi dây treo vật có giá trị
\choice
{$51,9\ N$}
{\True $30\ N$}
{$60\ N$}
{$42,4\ N$}

\haidong{10}
\loigiai{Chọn B

}\haidong{8} \end{ex}
\begin{ex}%book
Một vật có trọng lượng $P=25\ N$ đứng cân bằng nhờ hai dây $OA$ làm với trần một góc $60^{\circ}$ và $OB$ nằm ngang. Độ lớn của lực căng $\vec{T}_{A}$, $\vec{T}_{B}$ của dây $OA, OB$.
\choice
{$T_{A}=25 \sqrt{3}\ N ;\ T_{B}=25\ N$}
{$T_{A}=\dfrac{25}{\sqrt{3}}\ N ;\ T_{B}=\dfrac{50}{\sqrt{3}}\ N$}
{$T_{A}=25\ N ;\ T_{B}=25 \sqrt{3}\ N$}
{\True $T_{A}=\dfrac{50}{\sqrt{3}}\ N ;\ T_{B}=\dfrac{25}{\sqrt{3}}\ N$}	
\haidong{4}
\loigiai{Chọn D		
}\haidong{8} \end{ex}
\begin{ex}%book
Một quả cam có khối lượng $200 \ g$ đặt ở nơi có gia tốc rơi tự do là $g=10 \ m / s^2$. Trọng lượng của quả cam là
\choice
{\True $2 \ N$}
{\True $20 \ N$}
{$200 \ N$}
{$2000 \ N$}
\loigiai{

Chọn A

$P=0,2.10=2 N$

}
\haidong{10} \end{ex}
\begin{ex}%book%[1P2K5-4][Loai1]
Trọng lượng của một vật có khối lượng $20\ kg$ khi vật ở mặt đất có $g = 10\ m/s^2$ là
\choice
{$0,2\ N$}
{\True $200\ N$}
{$2\ N$}
{$20\ N$}
\loigiai{
Trọng lượng của vật là: $P = mg = 20.10 = 200\ N$.
}
\haidong{10} \end{ex}
\begin{ex}%book
Gia tốc rơi tự do trên Trái Đất là $g=9,8\ m/s^2$, gia tốc trên Mặt Trăng nhỏ hơn gia tốc trên Trái Đất $6$ lần. Nếu bạn Dũng có khối lượng 60 kg trên Trái Đất được lên Mặt Trăng du hành thì trọng lượng của bạn lúc này là
\choice
{$588 \ N$}
{\True $98 \ N$}
{$3528 \ N$}
{$600 \ N$}
\loigiai{

Chọn B

$P_{M T}=m \cdot g_{M T}=60 \cdot \dfrac{g_{T \ominus}}{6}=98 N$

}
\haidong{10} \end{ex}
\begin{ex}%book
Một ngọn đèn có khối lượng $m=1 \ kg$ được treo dưới trần nhà bằng một sợi dây. Lấy $g=9,8\ m/s^2$. Dây chỉ chịu lực căng lớn nhất là $8 \ N$. Nếu treo ngọn đèn này vào một đầu dây thì
\choice
{lực căng sợi dây là $9 \ N$ và dây sẽ bị đứt}
{\True lực căng sợi dây là $9,8 \ N$ và dây sẽ bị đứt}
{lực căng sợi dây là $9,8 \ N$ và dây sẽ không bị đứt}
{lực căng sợi dây là $4,9 \ N$ và dây sẽ không bị đứt}
\loigiai{
Chọn B

$T=P=m g=9,8 N>T_{\max}$ dây bị đứt

}
\haidong{10} \end{ex}
\begin{ex}%book%SBT KNTT 17.3
Biết gia tốc rơi tự do ở đỉnh và ở chân một ngọn núi lần lượt là $9,809\ m/s^2$ và $9,810\ m/s^2$. Tỉ số trọng lượng của vật ở đỉnh núi và chân núi là
\choice
{\True $0,9999$}
{$1,0001$}
{$9,8095$}
{$0,0005$}
\loigiai{
\begin{itemize}
\item Tỉ số trọng lượng: $\dfrac{P_1}{P_2}=\dfrac{g_1}{g_2}=\dfrac{9,809}{9,810}=0,9999$
\end{itemize}
}
\haidong{10} \end{ex}
\begin{ex}%book
Một người thả diều, dây diều được kéo lên cao tạo với mặt đất một góc $60^\circ$. Dùng lực kế để đo lực căng dây, người đó thấy rằng cánh diều đang kéo với lực $30 \ N$. Thành phần nằm ngang của lực kéo do diều sinh ra có độ lớn là
\choice
{$60 \ N$}
{$180 \ N$}
{\True $15 \ N$}
{$30 \ N$}
\loigiai{

Chọn C

$T_1=T \cos \alpha=30 \cos 60=15 N$.

}
\haidong{10} %<MyLT1>
\end{ex}
\begin{ex}%book
Các công nhân muốn treo một tấm biển nặng $40 \ kg$ vào trần nhà bằng dây cáp, nhưng dây cáp mà họ có sẵn chỉ có thể chịu lực tối đa 10 N. Lấy $g=10\ m/s^2$. Nếu họ định treo cáp ở một góc $45^\circ$, thì số lượng cáp ít nhất đủ để đỡ biển báo mà không bị đứt cáp là

\choice
{$20$ dây cáp}
{$80$ dây cáp}
{\True $60$ dây cáp}
{$40$ dây cáp}
\loigiai{
$n T_{\max} \sin 45^\circ=P \Rightarrow n=60$

}
\haidong{10} \end{ex}




















\subsubsection*{PHẦN 2. Thí sinh trả lời từ câu 1 đến câu 4. Trong mỗi ý a), b), c), d) ở mỗi câu, thí sinh chọn đúng hoặc sai.}
\setcounter{ex}{0}
\begin{ex}
Một chất điểm chịu tác dụng của hai lực $F_1=6\ N$ và $F_2=8\ N$, hợp với nhau một góc $\alpha$. Độ lớn của hợp lực là $F$.
\choiceTF[t]
{\True $\alpha=0^\circ$, thì $F=14\ N$}
{\True $\alpha=180^\circ$, thì $F=2\ N$}
{\True $\alpha=90^\circ$, thì $F=10\  N$}
{Nếu $\alpha=150^\circ$, thì $F \approx 12,44\ N$}
\loigiai{
\begin{itemchoice}
\itemch $\alpha=0^\circ \rightarrow F=F_1+F_2=6+8=14N$.
$\Rightarrow$ Đúng
\itemch $\alpha=180^\circ \rightarrow F=\left|F_1-F_2\right|=2N$.
$\Rightarrow$ Đúng
\itemch $\alpha=90^\circ \rightarrow F=\sqrt{F_1^2+F_2^2}=10N$.
$\Rightarrow$ Đúng
\itemch $\alpha=150^\circ \rightarrow F=\sqrt{F_1^2+F_2^2+2 F_1 \cdot F_2 \cos 150^\circ} \approx 4,11N$
$\Rightarrow$ Sai
\end{itemchoice}
}
\end{ex}	
\begin{ex}%Tailieuchuan
\immini[thm]{Một cậu bé đang kéo thùng hàng với vận tốc không đổi trên mặt đất bằng sợi dây hợp với phương ngang một góc $30^\circ$ như hình vẽ. Biết độ lớn lực kéo của cậu bé tác dụng dọc theo sợi dây là $12\ N$.}{
\includegraphics[width=4cm]{images/2025_04_24_b74eebea538181fc07eeg-18}}
\choiceTF[t]
{\True Chuyển động của thùng hàng là chuyển động thẳng đều} 
{Lực tổng hợp tác dụng lên thùng hàng không đổi và khác không} 
{\True Theo phương song song với mặt đất, lực kéo thành phần có độ lớn $6 \sqrt{3}\ N$} 
{\True Theo phương vuông góc với mặt đất, lực kéo thành phần có độ lớn $6\ N$} 
\loigiai{}
\end{ex}
\begin{ex}%Tailieuchuan
\immini[thm]{Một ô tô chịu một lực $F_1=400 \ N$ hướng về phía trước và một lực $F_2=300 \ N$ hướng về phía sau như hình vẽ.}{
\includegraphics[width=4cm]{images/2025_04_24_b74eebea538181fc07eeg-12}}
\choiceTF[t]
{\True $\vec{F}_1$ và $\vec{F}_2$ là hai lực cùng phương ngược chiều} 
{\True Ô tô chạy về phía trước} 
{Độ lớn hợp lực tác dụng lên ô tô được tính theo biểu thức: $F=F_1+F_2$} 
{Độ lớn hợp lực tác dụng lên ô tô là $700\ N$} 
\loigiai{}
\end{ex}
\begin{ex}%Tailieuchuan
\immini[thm]{Một bóng đèn được treo tại chính giữa một dây nằm ngang làm dây bị võng xuống. Biết trọng lượng của đèn là $100\ N$ và góc giữa hai nhánh dây là $150^\circ$.}{\begin{tikzpicture}[scale=1,thick,>=stealth',draw=Mapcolor]
\tkzDefPoints{0/0/a}
\tkzDefShiftPoint[a](-10.5:3){b}
\tkzDefShiftPoint[b](10:3){c}
\tkzDefShiftPoint[b](-90:1){m}
\draw[dashed] (-0.5,0)--(6.5,0);
\tkzDrawSegments(a,b b,c b,m)
\tkzDrawPoints[size=3,fill=white](b)

\draw[fill=Mapcolor!30] (m) ellipse (0.2cm and 0.3cm);
\node  [shape=rectangle,draw=black,fill=Mapcolor!30,minimum height=0.05cm,minimum width=0.1cm] at (2.95,-1.2) {};
\draw (2.6,-0.5) .. controls (2.8,-0.4) and (3.1,-0.4) .. (3.3,-0.5);
\node at (3,-0.2) {$150^\circ$};
\end{tikzpicture}}
\choiceTF[t]
{\True Lực tác dụng lên bóng đèn gồm 3 lực: $\vec{F}_1$, $\vec{\ F}_2$, $\vec{P}$ như hình vẽ} 
{\True Khi bóng đèn cân bằng: $\vec{P}+\vec{F}_1+\vec{F}_2=\vec{0}$} 
{\True Hợp lực của hai lực căng dây cân bằng với trọng lực} 
{Lực căng dây $F_1=F_2=170 \ N$} 
\loigiai{

}
\end{ex}

\subsubsection*{PHẦN 3. Thí sinh trả lời từ câu 1 đến câu 6.}
\setcounter{ex}{0} 
\begin{ex}%DoanhNamDinh
Một chất điểm chịu tác dụng của hai lực $F_1=9 \ N$ và $F_2=12 \ N$. Hợp lực có độ lớn $15\ N$, góc hợp bởi $\left(\vec{F_1}, \vec{F_2}\right)$ bằng bao nhiêu độ?
\shortans[oly]{90}

\loigiai{
Nhận thấy $F^2=F_1^2+F_2^2 \rightarrow \vec{F_1}$ vuông góc với $\vec{F_2}$, góc hợp bỡi $\left(\vec{F_1}, \vec{F_2}\right)=90^\circ$.
}
\end{ex}
\begin{ex}%DoanhNamDinh
Một chất điểm chịu tác dụng của hai lực $F_1=10 \ N$ và $F_2$, hợp với nhau một góc $150^\circ$. Độ lớn của hợp lực bằng $F=5\ N$. Giá trị của $F_2$ bằng bao nhiêu $N$ (làm tròn kết quả đến chữ số hàng phần trăm)?
\shortans[oly]{8,66}
\loigiai{
$F^2=F_1^2+F_2^2+2 F_1 F_2 \cos \alpha \rightarrow 25=100+F_2^2-10 \sqrt{3} F_2 \rightarrow F_2=5 \sqrt{3} N \approx 8,66\ N$.
}
\end{ex}
\begin{ex}%DoanhNamDinh
Một chất điểm chịu tác dụng của hai lực $F_1=10 \ N, \ F_2=5 \ N$. Hợp lực của hai lực $F=5 \sqrt{3}\ N$. Góc hợp bởi $\left(\vec{F_1}, \vec{F}\right)$ bằng bao nhiêu độ?
\shortans[oly]{30}
\loigiai{
$F^2=F_1^2+F_2^2+2 F_1 F_2 \cos \alpha \rightarrow \cos \alpha=\dfrac{F^2-F_1^2-F_2^2}{2 F_1 F_2}=-0,5 \rightarrow \alpha=120^{o}$
Nhận thấy: $F_1^2=F^2+F_2^2 \rightarrow \vec{F}$ vuông góc với $\vec{F_2}$. Vậy góc giữa $\left(\vec{F_1}, \vec{F}\right)$ bằng $30^{o}$.

}
\end{ex}


\begin{ex}%DoanhNamDinh
Một lít dầu có khối lượng $909\ g$. Lấy $g=9,8 \ m/s^2$. Một lượng dầu có trọng lượng $1500\ N$ sẽ đong được bao nhiêu lít dầu (làm tròn kết quả đến chữ số hàng đơn vị)?
\shortans[oly]{168}
\loigiai{

\begin{itemize}
\item Khối lượng dầu là $m=\dfrac{P}{g}=\dfrac{7500}{49} \ kg=\dfrac{7500000}{49} \ g$
\item Số lít dầu là $V=\dfrac{m}{909}=168,38=0,168 m^3$
\end{itemize}

}
\end{ex}


\begin{ex}%DoanhNamDinh
\immini[thm]{Một vật có trọng lượng $P=25 \ N$ được treo trên hai sợi dây mảnh $CA$, $CB$ (các dây nhẹ và không dãn) gắn trên trần nhà và tường như hình vẽ. Cho biết $\alpha=30^\circ$. Lực căng dây treo $CA$ bằng bao nhiêu niutơn?
}{\begin{tikzpicture}[scale=1,thick,>=stealth',draw=Mapcolor]
\draw (-0.5,0)-- (3.5,0)-- (3.5,-4);
\fill[pattern=north east lines] (-0.5,0)-- (3.5,0)-- (3.5,-4)-- (3.75,-4)-- (3.75,0.25)-- (-0.5,0.25);
\draw (1.5,-2)[very thick,fill=white] circle(0.11) coordinate (o)node[black, left=5pt]{{C}};
\draw (o)--++(0:2) coordinate (a)node[right=6pt]{{B}};
\draw (o)--++(120:2.3) coordinate (b) node[below left]{{A}};
\draw (o)--++(270:1.5) coordinate (m);
\shade[ball color=Mapcolor] (m) circle (1.5ex);
\node[white] at (m) {{m}};
\fill (1.5,-2)[white] circle(0.11);
\node at (1,-0.5) {$\alpha$};
\end{tikzpicture}}
\shortans[oly]{50}
\loigiai{
\includegraphics[width=4cm]{images/2025_06_04_af03f1d88c157fe9bf98g-182}

\begin{itemize}
\item Vật nằm cân bằng nên: $\vec{T_{A}}+\vec{T_{B}}+\vec{P}=\vec{0}$ (*)
\item Theo trục Oy, từ phương trình $\left({}^{*}\right)$ ta có: $T_{Al}-P=0 \Rightarrow \dfrac{\ T_{A}}{2}=P \Rightarrow T_{A}=50 \ N$
\end{itemize}

}
\end{ex}
\begin{ex}%DoanhNamDinh
\immini[thm]{Một chiếc áo có khối lượng $500\ g$ được treo vào điểm chính giữa của một sợi dây căng ngang, dây bị chùng xuống, hai nửa sợi dây có chiều dài như nhau và hợp với nhau một góc $120^\circ$ như hình vẽ. Lấy $g=9,8 \ m/s^2$, lực căng dây của mỗi nửa sợi dây bằng bao nhiêu niutơn?
}{\begin{tikzpicture}[draw=Mapcolor,scale=0.75]
\tkzInit[ymin=-2.5,ymax=2,xmin=-5.5,xmax=3.5]\tkzClip
\tkzDefPoints{0/0/o}
\tkzDefShiftPoint[o](0:2cm){o1}
\tkzDefShiftPoint[o](180:3cm){o2}
\tkzDefShiftPoint[o](26:3cm){b}
\tkzDefShiftPoint[o](162:5cm){a}
\tkzDefShiftPoint[a](-90:4cm){d}
\tkzDefShiftPoint[d](180:1cm){e}
\tkzDefShiftPoint[e](0:9.5cm){f}
\tkzDefPointBy[projection=onto e--f](b) \tkzGetPoint{c}
\tkzDrawSegments[line width=4pt,Mapcolor](a,d b,c)

\draw[orange] (-4.8,1.2) -- (-1,0.5) -- (2.6,1.2);
\draw[Mapcolor, very thick] (-1.1,0.4) .. controls (-1.1,0.7) and (-0.8,0.7) .. (-1,0.2);
\draw[Mapcolor, very thick] (-1,0.2) -- (-1.4,0.1);
\draw[Mapcolor, very thick] (-1,0.2) -- (-0.6,0.1);
\draw[Mapcolor, very thick] (-1.4,0.1) .. controls (-1.6,-0.2) and (-1.66,-0.5) .. (-1.66,-0.7) .. controls (-1.6,-0.7) and (-1.5,-0.7) .. (-1.44,-0.7) .. controls (-1.4,-0.5) and (-1.4,-0.3) .. (-1.27,-0.1) .. controls (-1.3,-0.3) and (-1.3,-0.5) .. (-1.3,-0.8) .. controls (-1.2,-0.8) and (-1.1,-0.8) .. (-1.04,-0.8) .. controls (-1,-0.5) and (-1,-0.3) .. (-1.1,0.17);
\draw[Mapcolor, very thick] (-0.6,0.1) .. controls (-0.4,-0.2) and (-0.34,-0.5) .. (-0.3,-0.7) .. controls (-0.4,-0.7) and (-0.47,-0.7) .. (-0.53,-0.7) .. controls (-0.57,-0.5) and (-0.6,-0.3) .. (-0.68,-0.1) .. controls (-0.7,-0.3) and (-0.7,-0.5) .. (-0.7,-0.8) .. controls (-0.8,-0.8) and (-0.9,-0.8) .. (-0.97,-0.8) .. controls (-0.95,-0.5) and (-0.9,-0.3) .. (-0.86,0.17);
\draw[] (-1.5,0.6) .. controls (-1.2,0.9) and (-0.8,0.9) .. (-0.5,0.6);
\node at (-1,1.1) {$120^\circ$};
\end{tikzpicture}}
\shortans[oly]{4,9}	
\loigiai{

\begin{itemize}
\item Vì chiếc áo đang nằm cân bằng nên: $\vec{T}_1+\vec{T}_2+\vec{P}=\vec{0}$
\item Điểm treo áo nằm chính giữa dây, do đó: $T_1=T_2=T$
\item Từ điều kiện cân bằng, ta có:
\includegraphics[width=4cm]{images/2025_06_04_af03f1d88c157fe9bf98g-182(1)}
$\Rightarrow T=\dfrac{m g}{2 \cdot \cos 60^\circ}=\dfrac{0,5 \cdot 9,8}{2 \cdot \cos 60^\circ}=4,9\ N$
\end{itemize}

}
\end{ex}







